\documentclass[10pt,a4paper]{amsart}
\usepackage[margin=19mm]{geometry}
\usepackage{amsmath,amssymb,mathtools}
\usepackage[scr=rsfs,cal=euler]{mathalfa}
\usepackage{xfrac}
\usepackage[unicode,hidelinks,bookmarks=false]{hyperref}
\newcommand{\ie}{i.e.\ }
\newcommand{\cf}{cf.\ }
\newcommand{\resp}{respectively\ }
\newcommand{\Subsection}{\subsection}
\providecommand{\nobreakdash}{\nobreak\hbox{-}}
\setlength{\emergencystretch}{2em}
\multlinegap0pt
\usepackage{amssymb}
\usepackage{mathtools}
\newtheorem{lem}{Lemma}
\DeclareMathOperator{\sign}{sgn}
\DeclareMathOperator{\spn}{span}
\title{On the free LAnKe on $3n-2$ generators: a theorem of Friedmann, Hanlon, Stanley and Wachs}
\author{Mihalis Maliakas, Dimitra-Dionysia Stergiopoulou}
\date{Source: arXiv:2401.09405v4 (CC BY 4.0)}
\begin{document}
\maketitle

\noindent\emph{Source and licence.} On the free LAnKe on $3n-2$ generators: a theorem of Friedmann, Hanlon, Stanley and Wachs. Version arXiv:2401.09405v4 (CC BY 4.0), distributed by arXiv under the Creative Commons Attribution 4.0 International licence, http://creativecommons.org/licenses/by/4.0/, as recorded in the arXiv licence metadata for this exact version and shown on the abstract page https://arxiv.org/abs/2401.09405v4. Full licence notice: \texttt{licenses/arxiv-2401.09405-CC-BY-4.0.txt}.\par\smallskip

\noindent\textit{Editorial scope.} Counted unit: the article's lem presentationlanke (source0.tex lines 848--849). The article's complete original proof of it is delivered with it (source0.tex lines 850--897, one \texttt{proof} environment of 48 lines). No mathematics is written, repaired or re-derived: the statement and the proof are the article's own lines, in the article's own notation, and no retained line is substituted or re-flowed.  Retained uncounted as the source-local context the proof needs: lines 761--763; lines 765--770; lines 778--784; lines 811--812; lines 832--836.  Nothing was omitted from any retained span, so the package carries no omission record.  No retained line contains a citation, so the wrapper carries no bibliography and no bibliography entry.  No retained line contains a figure environment, an image-inclusion command or drawing code, so no image is delivered and the package declares no asset dependency.  Editorial formatting only, in addition: an AMS \texttt{amsart} preamble that restates the article's own declarations and macros used by the retained lines, namely the environments \texttt{lem} on separate counters, together with the standard packages those lines need.  The retained environments print in this document as Lemma 1 in order of appearance, whereas the article prints the same items with the numbering of its own full document.  The complete native source of the article as distributed in the arXiv e-print for this exact version is bundled unchanged as \texttt{sources/153110/source0.tex}.
	\begin{align}
		\label{51}\textup{(G1)}&- \sign(\sigma)[[[{x_{\sigma(1)}},\dots ,{x_{\sigma(n)}} ],{y_{1}},\dots, y_{n-1}], z_{1},\dots ,z_{n-1}], \ \sigma \in \mathfrak{S}_n, \\\label{52}\textup{(G1)}&-\sign(\tau)[[[{x_{1}},\dots ,{x_{n}} ],{y_{\tau(1)}},\dots, y_{\tau(n-1)}], z_{1},\dots ,z_{n-1}], \ \tau \in \mathfrak{S}_{n-1},\\\label{53}\textup{(G1)}&-\sign(\tau)[[[{x_{1}},\dots ,{x_{n}} ],{y_{1}},\dots, y_{n-1}], z_{\tau(1)},\dots ,z_{\tau(n-1)}], \ \tau \in \mathfrak{S}_{n-1},
	\end{align}

	\begin{align}
		\label{54}\textup{(G2)}&-\sign(\sigma)[[{x_{\sigma(1)}},\dots ,{x_{\sigma(n)}} ],[{y_{1}},\dots, y_{n}],  z_{1},\dots ,z_{n-2}  ],  \ \sigma \in \mathfrak{S}_n,
		\\\label{55}\textup{(G2)}&-\sign(\sigma)[[{x_1},\dots ,{x_n} ],[{y_{\sigma(1)}},\dots, y_{\sigma(n)}], z_{1},\dots ,z_{n-2}  ], \ \sigma \in \mathfrak{S}_n,\\
		\label{56}\textup{(G2)}&-\sign(\tau)[[{x_1},\dots ,{x_n} ],[{y_{1}},\dots, y_{n}], z_{\tau(1)},\dots ,z_{\tau(n-2)}  ], \ \tau \in \mathfrak{S}_{n-2},
		\\\label{57}\textup{(G2)}&+ [[{y_1},\dots ,{y_n} ],[{x_{1}},\dots, x_{n}], z_{1},\dots ,z_{n-2} ].
	\end{align}

\begin{lem}\label{genrel1}
	Let $\spn (R1,R2,R3)$ be the subspace of $W$ spanned by the elements $\textup{(R1), (R2), (R3)}$  defined as follows   \begin{align}\tag{R1} \textup{(G1)}&- \sum_{i=1}^{n}(-1)^{i-1}[[[x_i,{y_{1}},\dots, y_{n-1}],{x_1},\dots ,\widehat{x_i}, \dots, x_n ],z_{1},\dots ,z_{n-1}],\\
		\tag{R2} \textup{(G1)}&- [[[{x_1},\dots ,{x_n} ],{z_{1}},\dots, z_{n-1}], y_{1},\dots ,y_{n-1}]
		\\\nonumber&-\sum_{i=1}^{n-1}(-1)^{i-1}[[{x_1},\dots ,{x_{n}}],[y_{i},z_1, \dots, z_{n-1} ],{y_{1}},\dots, \widehat{y_i},\dots ,y_{n-1}],\\\tag{R3} \textup{(G2)}&-\sum_{i=1}^{n}(-1)^{i}[[[y_1,\dots,y_n],x_i , z_1,\dots, z_{n-2}],x_1,\dots,\widehat{x_i},\dots, x_n],
	\end{align}
for all $x_i, y_j, z_k \in [m]$. Then the $\mathfrak{S_m}$-modules  $\text{Lie}_n(m)$ and $W/\spn (R1,R2,R3)$ are isomorphic. 
\end{lem}

\begin{lem}\label{liebarsur} The inclusion map $W(1) \subseteq W$ induces a surjective map of $\mathfrak{S}_m$-modules \[\overline{\textup{Lie}}_n(m) \to \textup{Lie}_n(m).\]
\end{lem}

\begin{lem}\label{finalrel}
In $W(1)$  we have the relations 
\begin{align*}&\sum_{i=1}^n(-1)^i[[[y_1,\dots,y_n],x_i,z_1,\dots, z_{n-2}],x_1,\dots, \widehat{x_i}, \dots, x_n]\\&-\sign(\sigma)\sum_{i=1}^n(-1)^i[[[y_1,\dots,y_n],x_{\sigma(i)},z_1,\dots, z_{n-2}],x_{\sigma(1)},\dots, \widehat{x_{\sigma(i)}}, \dots, x_{\sigma(n)}],
\end{align*} where $\sigma \in \mathfrak{S}_n$. Thus we have the corresponding relations in $\overline{\textup{Lie}}_n(m)$.
\end{lem}

\begin{lem}\label{presentationlanke} The inclusion map $W(1) \subseteq W$ induces an isomorphism of $\mathfrak{S}_m$-modules \[\overline{\textup{Lie}}_n(m) \to \textup{Lie}_n(m).\]
\end{lem}

\begin{proof}  From Lemma \ref{liebarsur} we know that the inclusion map $W(1) \subseteq W$ induces a surjective linear map  $ \overline{\text{Lie}}_n(m) \to \text{Lie}_n(m).$ 
	
	In the argument given below, we show that there exists a surjective linear map $\text{Lie}_n(m) \to \overline{\text{Lie}}_n(m)$. Since these spaces are finite dimensional, the surjective linear map $\overline{\text{Lie}}_n(m) \to \text{Lie}_n(m)$  is an isomorphism as desired.
	
	Consider the map \[ \Theta:W \to \overline{\text{Lie}}_n(m)\]  \begin{itemize}
		\item that is the identity map on (G1), and
		\item sends a generator $[[{x_1},\dots ,{x_n} ],[{y_{1}},\dots, y_{n}], z_{1},\dots ,z_{n-2}  ]$ of type (G2) to 
		\begin{equation}\label{ThetaG2}\sum_{i=1}^{n}(-1)^{i}[[[y_1,\dots,y_n],x_i , z_1,\dots, z_{n-2}],x_1,\dots,\widehat{x_{i}},\dots, x_n].\end{equation}
	\end{itemize}
	We need to show that $\Theta$ is well defined, i.e. sends the relations  of  $W$ to relations of $\overline{\text{Lie}}_n(m)$ .
	
	Since $\Theta$ is the identity map on (G1), it is clear that the relations (\ref{51}) - (\ref{53})  are sent to the corresponding relations of $\overline{\text{Lie}}_n(m)$.
	
	Next we consider the relations (\ref{54}) - (\ref{57}) of $W$. From the definition of $\Theta$, it follows that the image of the relation (\ref{54}) is equal to 
	\begin{align*}&\sum_{i=1}^n(-1)^i[[[y_1,\dots,y_n],x_i,z_1,\dots, z_{n-2}],x_1,\dots, \widehat{x_i}, \dots, x_n]\\&-\sign(\sigma)\sum_{i=1}^n(-1)^i[[[y_1,\dots,y_n],x_{\sigma(i)},z_1,\dots, z_{n-2}],x_{\sigma(1)},\dots, \widehat{x_{\sigma(i)}}, \dots, x_{\sigma(n)}].
	\end{align*}
	By Lemma \ref{finalrel} this is equal to $0$.
	
		In order to show that relation (\ref{55}) of  $W$ is mapped under $\Theta$ to a relation of $\overline{\text{Lie}}_n(m)$, it suffices by the definition given in (\ref{ThetaG2}) to show the following identity in $\overline{\text{Lie}}_n(m)$
	\begin{align}\label{ThetaG21} &\sum_{i=1}^{n}(-1)^{i}[[[y_1,\dots,y_n],x_i , z_1,\dots, z_{n-2}],x_1,\dots,\widehat{x_i},\dots, x_n] \\\nonumber=&\sign(\sigma)\sum_{i=1}^{n}(-1)^{i}[[[y_{\sigma(1)},\dots,y_{\sigma(n)}],x_i , z_1,\dots, z_{n-2}],x_1,\dots,\widehat{x_i},\dots, x_n] \end{align}
	for any permutation $\sigma \in \mathfrak{S}_n$. From relation (\ref{51}) of $W(1)$ applied to   \[[[[y_1,\dots,y_n],x_i , z_1,\dots, z_{n-2}],x_1,\dots,\widehat{x_i},\dots, x_n] \]
	we have
	\begin{align*}&[[[y_1,\dots,y_n],x_i , z_1,\dots, z_{n-2}],x_1,\dots,\widehat{x_i},\dots, x_n] \\\nonumber=& \sign(\sigma) [[[y_{\sigma(1)},\dots,y_{\sigma(n)}],x_i , z_1,\dots, z_{n-2}],x_1,\dots,\widehat{x_i},\dots, x_n]. \end{align*}
	Taking the alternating sum with respect to $i=1,\dots, n$ we obtain eq. (\ref{ThetaG21}).
	
		Similarly to the previous case, in order to show that relation (\ref{56}) of  $W$ is mapped under $\Theta$ to a relation of $\overline{\text{Lie}}_n(m)$, it suffices to show the following identity in $\overline{\text{Lie}}_n(m)$
	\begin{align}\label{ThetaG22} &\sum_{i=1}^{n}(-1)^{i}[[[y_1,\dots,y_n],x_i , z_1,\dots, z_{n-2}],x_1,\dots,\widehat{x_i},\dots, x_n] \\\nonumber=& \sign(\tau)\sum_{i=1}^{n}(-1)^{i}[[[y_{1},\dots,y_{n}],x_i , z_{\tau(1)},\dots, z_{\tau(n-2)}],x_1,\dots,\widehat{x_i},\dots, x_n] \end{align}
	for any permutation $\tau \in \mathfrak{S}_{n-2}$. From relation (\ref{52}) of $W(1)$ applied to   \[[[[y_1,\dots,y_n],x_i , z_1,\dots, z_{n-2}],x_1,\dots,\widehat{x_i},\dots, x_n] \]
	we have
	\begin{align*}&[[[y_1,\dots,y_n],x_i , z_1,\dots, z_{n-2}],x_1,\dots,\widehat{x_i},\dots, x_n] \\\nonumber=&\sign(\tau)[[[y_{1},\dots,y_{n}],x_i , z_{\tau(1)},\dots, z_{\tau(n-2)}],x_1,\dots,\widehat{x_i},\dots, x_n], \end{align*}
	where the element $x_i$ remains fixed. Hence (\ref{ThetaG22}) follows.
	
	From the definition of $\Theta$ (see (\ref{ThetaG2})), it follows that the image of the relation (\ref{57}) is equal to the relation (R5) of $\overline{\text{Lie}}_n(m)$. 
	
	Thus far we have shown that $\Theta:W \to \overline{\text{Lie}}_n(m)$ is a well defined map. Since $\overline{\text{Lie}}_n(m)$ is a quotient of $W(1)$ and the (G1) generate $W(1)$, it is clear from the definition of $\Theta$ that $\Theta$ is surjective.
	
	Finally, from the definition of $\Theta$ it is clear that the image under $\Theta$
	\begin{itemize}
		\item of (R1) of Lemma \ref{genrel1} is equal to (R1) in $\overline{\textup{Lie}}_n(m)$,
		\item of (R2) of Lemma \ref{genrel1} is equal to (R4) in $\overline{\textup{Lie}}_n(m)$, and
		\item of (R3) of Lemma \ref{genrel1} is equal to zero.
	\end{itemize}
From Lemma \ref{genrel1} it follows that $\Theta$ induces a surjective linear map $\text{Lie}_n(m) \to \overline{\text{Lie}}_n(m)$. The proof is complete.
	
	

	
 	\end{proof}

\end{document}
